\thesisappendix

\chapter{EMAT通信协议规范}

\section{数据包格式}

EMAT探头与主控之间采用自定义二进制通信协议，数据包格式如下：

\begin{center}
\begin{tabular}{|c|c|c|c|c|c|}
\hline
\textbf{帧头} & \textbf{保留} & \textbf{保留} & \textbf{功能码} & \textbf{长度(H)} & \textbf{长度(L)} \\
\hline
0xAB & 0x00 & 0x01 & 见表 & MSB & LSB \\
\hline
\end{tabular}
\end{center}

载荷字段紧随长度字段之后，最后为1字节CRC-8校验（多项式0x07，初始值0x00）。

\section{功能码定义}

\begin{table}[htbp]
\centering
\caption{功能码与对应操作}
\label{tab:func_codes}
\begin{tabular}{clp{6cm}}
\toprule
\textbf{功能码} & \textbf{命令} & \textbf{说明} \\
\midrule
0x00 & 读取厚度 & 请求当前厚度测量值（预留）\\
0x01 & 读取波形 & 波形数据分4块传输，每块约8185采样点\\
0x03 & 设置参数 & 写入激励频率、平均次数和声速\\
0x04 & 获取参数 & 读取当前采集参数\\
\bottomrule
\end{tabular}
\end{table}

\chapter{主要符号表}

\begin{table}[htbp]
\centering
\caption{本文主要符号及其含义}
\begin{tabular}{cl}
\toprule
\textbf{符号} & \textbf{含义} \\
\midrule
$\mathbf{L}$ & 洛伦兹体力（N/m$^3$）\\
$\mathbf{J}$ & 电流密度（A/m$^2$）\\
$\mathbf{B}$ & 磁通密度（T）\\
$\sigma$ & 材料电导率（S/m）\\
$\rho$ & 材料密度（kg/m$^3$）\\
$\lambda, \mu$ & Lam\'{e}常数（Pa）\\
$c_L, c_T$ & 纵波/横波速度（m/s）\\
$s(t)$ & 时域回波信号\\
$h$ & 升距（m）\\
$\hat{p}_t$ & 时刻$t$的接触概率估计\\
$\mathbf{P}$ & 物理约束矩阵\\
$g_m$ & 模态$m$的可靠性门控权重\\
\bottomrule
\end{tabular}
\end{table}